% BHU PYQ -- Auto-generated & error-corrected
\documentclass[11pt,a4paper]{article}

\usepackage[a4paper,top=2.5cm,bottom=2.5cm,left=2.8cm,right=2.8cm,headheight=14pt]{geometry}
\usepackage{amsmath,amssymb}
\usepackage[T1]{fontenc}
\usepackage[utf8]{inputenc}
\usepackage{lmodern}
\usepackage{microtype}
\usepackage{fancyhdr}
\usepackage{enumitem}
\usepackage{booktabs}
\usepackage{hyperref}
\usepackage{lastpage}
\usepackage{parskip}

\hypersetup{colorlinks=true,urlcolor=black,linkcolor=black,
  pdftitle={STA/STB-504: Operations Research -- BHU PYQ},pdfauthor={Banaras Hindu University}}

\pagestyle{fancy}\fancyhf{}
\renewcommand{\headrulewidth}{0.4pt}\renewcommand{\footrulewidth}{0.4pt}
\fancyhead[L]{\small   \textit{Banaras Hindu University}}
\fancyhead[R]{\small   \textit{STA/STB-504: Operations Research}}
\fancyfoot[C]{\small Page  \thepage\ of \pageref{LastPage}}


\newlist{parts}{enumerate}{2}
\setlist[parts,1]{label=(\alph*),leftmargin=*,itemsep=5pt,topsep=3pt}
\setlist[parts,2]{label=(\roman*),leftmargin=*,itemsep=3pt,topsep=2pt}

\newcommand{\pts}[1]{\hfill   \textit{[#1]}}


\begin{document}


\begin{center}
  {\large\bfseries BANARAS HINDU UNIVERSITY}\\[2pt]
  {
\normalsize B.A./B.Sc. (Hons.) Semester V Examination, 2022-23}\\[6pt]
  \rule{0.8\linewidth}{0.5pt}\\[6pt]
  {\large\bfseries Paper No. STA/STB-504: Operations Research}\\[4pt]
  {
\normalsize Time: 3 Hours\hfill Maximum Marks: 70}
\end{center}
\rule{\linewidth}{0.4pt}
\smallskip



\noindent   \textit{Instructions: Answer any   \textbf{five} questions. All questions carry equal marks. Symbols have their usual meanings.}
\bigskip

\medskip


\noindent   \textbf{Question 1.}

\begin{parts}
  \item Define a two-person zero-sum game. When is a game said to be strictly determinable? Solve the following $2  \times 2$ game algebraically:
  \[
     \text{Player A's payoff matrix:} \quad
    
\begin{pmatrix} 4 & 2 \\ 1 & 3 \end{pmatrix}
  \]\pts{7}
  \item Solve the following game using the matrix oddment (graphical or algebraic) method:
  \[
     \text{Player A's payoff matrix:} \quad
    
\begin{pmatrix} 1 & 3 & 2 \\ 4 & 2 & 5 \end{pmatrix}
  \]\pts{7}
\end{parts}

\medskip\rule{\linewidth}{0.2pt}

\medskip


\noindent   \textbf{Question 2.}

\begin{parts}
  \item Use the graphical method to convert the following game into a $2 \times 2$ game and obtain its solution:
  \[
     \text{Player A's payoff matrix:} \quad
    
\begin{pmatrix} 1 & 3 & 11 \\ 8 & 5 & 2 \end{pmatrix}
  \]\pts{7}
  \item Differentiate between mixed strategy and pure strategy in game theory. Let $(a_{ij})$ be the $m  \times n$ payoff matrix for a two-person zero-sum game. If $\underline{v}$ denotes the maximin value and $ar{v}$ denotes the minimax value of the game, then show that $\underline{v} \leq ar{v}$. \pts{7}
\end{parts}

\medskip\rule{\linewidth}{0.2pt}

\medskip


\noindent   \textbf{Question 3.}

\begin{parts}
  \item Given that $x_1 = 1$, $x_2 = 1$, $x_3 = 1$ is a feasible solution to the system of equations $x_1 + x_2 + x_3 = 3$, $2x_1 + x_2 = 3$. Reduce the given feasible solution to a basic feasible solution. \pts{7}
  \item When is an LPP said to have an infeasible or multiple optimal solution? Comment on the solutions of the following LPP:
  
\begin{parts}
    \item Maximise $Z = 3x_1 - x_2$ subject to: $-x_1 + x_2 \geq 0$, $x_1 + 2x_2 \leq -1$, $x_1, x_2 \geq 0$.
    \item Maximise $Z = 15x_1 + 30x_2$ subject to: $x_1 + 4x_2 \leq 60$, $x_1 + 2x_2 \leq 60$, $x_1, x_2 \geq 0$.
  \end{parts}\pts{7}
\end{parts}

\medskip\rule{\linewidth}{0.2pt}

\medskip


\noindent   \textbf{Question 4.}

Use the simplex procedure to solve the following LPP:

\begin{align*}
   \text{Maximise} \quad & Z = 2x_1 + x_2 \\
   \text{subject to} \quad & x_1 + 2x_2 \leq 10 \\
  & x_1 + x_2 \leq 6 \\
  & x_1 - x_2 \leq 2 \\
  & x_1 - 2x_2 \leq 1 \\
  & x_1, x_2 \geq 0
\end{align*}
Also state what the dual variable values are if we convert to the dual. \pts{14}

\medskip\rule{\linewidth}{0.2pt}

\medskip


\noindent   \textbf{Question 5.}

Use the two-phase method to solve the following LPP:

\begin{align*}
   \text{Maximise} \quad & Z = 3x_1 - x_2 \\
   \text{subject to} \quad & x_1 + x_2 \geq 2 \\
  & x_1 + 3x_2 \geq 5 \\
  & x_1 \leq 4 \\
  & x_1, x_2 \geq 0
\end{align*}\pts{14}

\medskip\rule{\linewidth}{0.2pt}

\medskip


\noindent   \textbf{Question 6.}

What are the types of inventory control models? Explain the different types of costs used in inventory control. Derive the EOQ (Economic Order Quantity) model with shortages. \pts{14}

\medskip\rule{\linewidth}{0.2pt}

\medskip


\noindent   \textbf{Question 7.}

Differentiate between a transportation problem and an assignment problem. The following table gives the profit (in Rs.) earned by assigning jobs 1, 2, 3, and 4 to contractors A, B, C, and D:


\begin{center}

\begin{tabular}{c|cccc}
 oprule
Job & A & B & C & D \\
\midrule
1 & 10 & 15 & 9 & 12 \\
2 & 11 & 14 & 10 & 13 \\
3 & 9 & 15 & 10 & 13 \\
4 & 8 & 14 & 12 & 11 \\
ottomrule
\end{tabular}
\end{center}

Convert the problem into an LPP and find the assignment that maximises the total profit. \pts{14}

\medskip\rule{\linewidth}{0.2pt}

\medskip


\noindent   \textbf{Question 8.}

Items are to be shipped from three sources to four destinations as per the following cost table. Using the MODI method, find the total minimum transportation cost:


\begin{center}

\begin{tabular}{c|cccc|c}
 oprule
 & $D_1$ & $D_2$ & $D_3$ & $D_4$ & Supply \\
\midrule
$S_1$ & 19 & 30 & 50 & 10 & 7 \\
$S_2$ & 70 & 30 & 40 & 60 & 9 \\
$S_3$ & 40 & 8 & 70 & 20 & 18 \\
\midrule
Demand & 5 & 8 & 7 & 14 & \\
ottomrule
\end{tabular}
\end{center}\pts{14}

\vfill
\rule{\linewidth}{0.4pt}

\begin{center}\small
  \href{https://bhu-exam.vercel.app/}{Click here to visit our website}\\[4pt]
  \href{https://me-aryan.vercel.app/}{Click here to visit our contributor}\\[4pt]
  \href{https://quashan-me.vercel.app/}{Click here to visit our contributor}
\end{center}

\end{document}